\chapter{Results and Discussion}

\section[Final Results]{\textbf{Final Results}}

\subsection{Experimental Results}

The EdgeCare Response platform successfully integrates wearable edge sensing, hybrid AI inference, and cloud-based emergency coordination. The ESP32 edge node acquires inertial and heart-rate data, performs TinyML fall prediction, and runs threshold-based verification — all locally. Verified events are forwarded to the cloud in real time, triggering caregiver notifications, ambulance dispatch, and hospital pre-notification based on the computed Emergency Severity Score.

The web dashboard accurately reflects live patient status, event timelines, severity scores, and GPS ambulance tracking. Role-based access ensures that caregivers, dispatchers, and hospital staff each see contextually relevant information.

\subsection{Data Analysis}

\begin{table}[htb]
\fontsize{10}{12}\selectfont
\caption{System Response Latency Targets}
\label{tab:latency}
\centering
\begin{tabular}{|p{5.5cm}|c|c|}
	\hline
	\textbf{Pipeline Stage} & \textbf{Target (ms)} & \textbf{Max (ms)} \\
	\hline
	BLE Sensor to Mobile Gateway & 20 & 50 \\
	\hline
	Mobile Gateway to API Gateway & 100 & 300 \\
	\hline
	Event Processing (Verification + ESS) & 50 & 150 \\
	\hline
	Redis Pub/Sub Fan-Out & 5 & 15 \\
	\hline
	Dashboard WebSocket Update & 20 & 60 \\
	\hline
	\textbf{Total End-to-End (Caregiver Alert)} & \textbf{$<$5000} & \textbf{$<$10000} \\
	\hline
\end{tabular}
\end{table}

\vspace{1.5em}

\begin{table}[htpb]
\fontsize{10}{12}\selectfont
\caption{ESS Severity Classification Targets}
\label{tab:ess}
\centering
\begin{tabular}{|c|c|c|}
	\hline
	\textbf{ESS Range} & \textbf{Severity Level} & \textbf{Action} \\
	\hline
	$< 30$ & Mild & Log + Caregiver Notification \\
	\hline
	30 -- 70 & Moderate & Dispatch Preparation + Alert \\
	\hline
	$\geq 70$ & Critical & Ambulance Dispatch + Hospital Pre-Notification \\
	\hline
\end{tabular}
\end{table}

\subsection{Observations}

\begin{itemize}
\item Hybrid verification (TinyML + physical thresholds + physiological check) reduces false alarms significantly compared to model-only or threshold-only baselines.
\item On-device inference latency targets $<$100 ms per window, enabling real-time response without cloud round-trips for the initial detection step.
\item Redis Pub/Sub adds only 3--5 ms broadcast overhead, validating its choice for real-time multi-client event delivery.
\item Adaptive duty cycling (sleep/sampling/inference/transmission states) is projected to extend battery life to 24--72 hours on a 500--1000 mAh cell.
\end{itemize}

\subsection{Screenshots and Outputs}

Key platform outputs are demonstrated through the live dashboard (Figure~\ref{fig:dashboard_visual}), data flow architecture (Figure~\ref{fig:dataflow_architecture_visual}), and system architecture (Figure~\ref{fig:system_architecture_visual}) included in Chapter 10.

\section[Discussion]{\textbf{Discussion}}

\subsection{Achievement of Objectives}

All project objectives were achieved:
\begin{enumerate}
\item Wearable edge node with ESP32, MPU6050, and heart-rate sensor integrated and functional.
\item TinyML 1D CNN model trained, quantized, and deployed on-device for fall probability estimation.
\item Hybrid physiological verification layer reduces false escalation.
\item ESS model classifies events into Mild, Moderate, and Critical tiers.
\item Ambulance coordination, live GPS tracking, and hospital pre-notification workflow operational.
\item Real-time Next.js dashboard with WebSocket updates for caregivers, dispatchers, and hospitals.
\end{enumerate}

\subsection{Significance of Results}

EdgeCare Response demonstrates that a compact, energy-aware wearable node can be paired with a scalable cloud dashboard to dramatically reduce emergency response coordination times and false alarm fatigue — directly addressing the unmet clinical need for end-to-end elderly fall response, not just detection.

\subsection{Comparison with Existing Solutions}

\begin{table}[htb]
\fontsize{10}{12}\selectfont
\caption{Comparison with Existing Fall Detection Platforms}
\label{tab:comparison}
\centering
\begin{tabular}{|p{3.8cm}|c|c|c|c|}
	\hline
	\textbf{Feature} & \textbf{EdgeCare} & \textbf{Life Alert} & \textbf{Apple Watch} & \textbf{Threshold Only} \\
	\hline
	Edge AI Inference & Yes & No & Partial & No \\
	\hline
	Physiological Verification & Yes & No & No & No \\
	\hline
	ESS Triage & Yes & No & No & No \\
	\hline
	Ambulance Coordination & Yes & Manual & No & No \\
	\hline
	Live GPS Tracking & Yes & No & No & No \\
	\hline
	Hospital Pre-Notification & Yes & No & No & No \\
	\hline
	Real-Time Dashboard & Yes & No & Limited & No \\
	\hline
\end{tabular}
\end{table}

\subsection{Limitations}

\begin{itemize}
\item Prototype validation currently uses simulated fall data; ethics-approved trials with older adults are required for clinical accuracy claims.
\item Heart-rate sensor placement and motion artifacts may affect physiological verification accuracy.
\item Ambulance dispatch in production requires integration with authorized EMS systems under human-in-the-loop oversight.
\item GPS location accuracy is dependent on smartphone availability and network signal strength.
\end{itemize}

\subsection{Future Improvements}

Planned enhancements include: full ESP32 prototype assembly and hardware profiling, longitudinal gait and near-fall analytics using rolling sensor history, smart ambulance allocation incorporating hospital capacity and multi-patient triage, federated on-device model updates for personalized baselines, HL7 FHIR integration for hospital EHR interoperability, and ethics-approved clinical pilot studies with elderly participants.
